\section{System Model}
\label{sec:model}

This section fixes notation for the rest of the paper and names the platform's artifacts.

\subsection{Network}

A Freenet \emph{peer} is an instance of the Freenet Core running on some host. Peers communicate over a UDP-based encrypted transport (\cref{sec:transport}). Each peer $p$ has a \emph{location} $\ell(p) \in [0, 1)$, computed as a cryptographic hash of the peer's external network address mapped into the unit interval. Locations live on a one-dimensional ring with wrap-around distance
\[
  d(x, y) = \min\bigl(|x - y|,\, 1 - |x - y|\bigr).
\]
A peer maintains a set of \emph{neighbor} connections; the connection set is bounded between configurable limits (typically 25 and 200) and is reorganized over time according to the topology rule of \cref{sec:smallworld}.

\subsection{Artifacts}

Three kinds of artifact live on top of the network.

\paragraph{Contracts.} A \emph{contract} $\contract$ is a tuple $(\mathrm{code}, \mathrm{params})$ where $\mathrm{code}$ is a WebAssembly module conforming to the platform's contract interface and $\mathrm{params}$ is an opaque byte string supplied at instantiation time. The \emph{key} of the contract is $k(\contract) = H(\mathrm{code} \,\|\, \mathrm{params})$ for a cryptographic hash $H$, and its location is $\ell(\contract) \in [0, 1)$ obtained from $k(\contract)$ in the same way as a peer location. Each contract has a \emph{state} $\sigma \in \statespace_\contract$, replicated on a subset of peers near $\ell(\contract)$ on the ring. The contract's code defines the structure of $\statespace_\contract$, the validity predicate, the merge operation, and the summary/delta protocol described in the next two sections. Contracts are public: their state is readable by any peer that can route to $\ell(\contract)$.

\paragraph{Delegates.} A \emph{delegate} is a WebAssembly module that runs on a single user's device. It holds private state (managed as secret storage by the local core) and responds to attributed messages from user interfaces, contracts, and other delegates (\cref{sec:delegates}). Delegates are local; they do not appear on the ring.

\paragraph{User interfaces.} A \emph{user interface} is a web frontend --- HTML, JavaScript, WebAssembly --- delivered by a contract whose state encodes the application bundle. User interfaces run in the browser, talk to the local core over a WebSocket API, and reach contracts and delegates indirectly through that core.

\subsection{Operations}
\label{sec:ops-summary}

A small set of operations connect these artifacts. We describe them informally here and revisit their lifecycles in \cref{sec:operations}:

\begin{description}[leftmargin=1.4em, style=nextline, itemsep=2pt]
  \item[\textsf{PUT}] Publish a contract: insert or update its state at the peers near $\ell(\contract)$.
  \item[\textsf{GET}] Retrieve a contract's current state by key.
  \item[\textsf{UPDATE}] Submit an update to a contract; the update is merged into existing replicas using the contract's merge function.
  \item[\textsf{SUBSCRIBE}] Receive notifications whenever the contract's state changes.
  \item[\textsf{CONNECT}] Establish a peering relationship with the network as a new node (\cref{sec:smallworld}).
\end{description}

\textsf{PUT}, \textsf{GET}, \textsf{UPDATE}, and \textsf{SUBSCRIBE} all route over the ring; their destination is $\ell(\contract)$. \textsf{CONNECT} is the bootstrap operation by which a fresh peer joins the topology.

\subsection{Trust Model}

Peers are not trusted to behave correctly. Messages are end-to-end authenticated by transport, but a peer along a routing path can drop, reorder, or fail to forward messages. Replication and adaptive routing (\cref{sec:adaptive}) provide robustness against unresponsive or slow peers; cryptographic validity predicates inside contracts (e.g.\ signature checks on updates) provide robustness against actively malicious peers attempting to inject invalid state. The platform does not assume a global trusted authority.

A user's own peer, and the delegates running on it, are trusted on behalf of that user. Delegates protect secrets against application bugs on the same device; they do not defend against a compromised operating system.
